
\documentclass[11pt,a4paper]{article}

\usepackage{fontspec}
\usepackage[french]{babel}
\usepackage{amsmath,amssymb,amsthm,mathtools}
\usepackage{geometry}
\usepackage{xcolor}
\usepackage{fancyhdr}
\usepackage{enumitem}
\usepackage{hyperref}
\usepackage{titlesec}
\usepackage{microtype}
\usepackage{booktabs}
\usepackage{array}
\usepackage{tabularx}
\usepackage{longtable}
\usepackage{colortbl}

\geometry{margin=2.2cm}
\setmainfont{DejaVu Serif}
\setsansfont{DejaVu Sans}
\setmonofont{DejaVu Sans Mono}

\definecolor{accent}{HTML}{1F4E79}
\definecolor{lightaccent}{HTML}{EAF2F8}
\definecolor{lightgray}{HTML}{F7F7F7}

\pagestyle{fancy}
\fancyhf{}
\lhead{\small Cours Deep Reinforcement Learning}
\chead{\small Classe : 4IA}
\rhead{\small Année : 2026}
\lfoot{\small Enseignant : Mourad Zéraï, PhD - ESPRIT School of Engineering}
\cfoot{}
\rfoot{\thepage}
\setlength{\headheight}{14pt}

\titleformat{\section}{\large\bfseries\color{accent}}{\thesection.}{0.5em}{}
\titleformat{\subsection}{\normalsize\bfseries\color{accent}}{\thesubsection.}{0.5em}{}
\titleformat{\subsubsection}{\normalsize\bfseries}{\thesubsubsection.}{0.5em}{}

\newtheorem*{definition}{Définition}
\newtheorem*{rappel}{Rappel}
\newtheorem*{hypothese}{Hypothèse}
\newtheorem*{propriete}{Propriété}
\newtheorem*{theoreme}{Théorème}
\newtheorem*{remarque}{Remarque}

\newcommand{\E}{\mathbb{E}}
\newcommand{\Pbb}{\mathbb{P}}
\newcommand{\R}{\mathbb{R}}
\newcommand{\argmax}{\operatorname*{argmax}}
\newcommand{\given}{\,|\,}

\begin{document}

\begin{center}
    {\LARGE\bfseries Dérivation et explication pas à pas de Q-Learning}\\[0.4em]
    {\large Note de cours rigoureuse, détaillée et appliquée à FrozenLake}
\end{center}

\vspace{0.6em}

\section{Objectif}

Le but de cette note est de dériver, pas à pas, l'algorithme de Q-Learning, de justifier chaque passage important et de montrer comment l'appliquer à FrozenLake.

\section{Cadre et point de départ}

On considère un MDP fini $(\mathcal S,\mathcal A,p,r,\gamma)$ avec $\gamma \in [0,1)$. On cherche la fonction état-action optimale
\[
q_*(s,a)=\max_{\pi} q_{\pi}(s,a).
\]

\begin{rappel}
La fonction optimale vérifie l'équation d'optimalité de Bellman
\[
q_*(s,a)=\sum_{s',r} p(s',r\given s,a)\left[r+\gamma \max_{a'} q_*(s',a')\right].
\]
\end{rappel}

Cette équation contient deux idées :
\begin{itemize}[leftmargin=1.6em]
    \item une \emph{espérance} sur la transition inconnue ;
    \item une \emph{maximisation} sur l'action future.
\end{itemize}
Q-Learning remplace l'espérance exacte par des échantillons successifs et conserve la maximisation.

\section{Passage du modèle exact à une mise à jour échantillonnée}

\subsection{Opérateur de Bellman optimal pour $q$}

On définit
\[
(T_*Q)(s,a)=\sum_{s',r} p(s',r\given s,a)\left[r+\gamma \max_{a'} Q(s',a')\right].
\]
Un point fixe de $T_*$ est une solution de l'équation d'optimalité de Bellman.

\subsection{Observation d'une seule transition}

Supposons qu'à l'instant $t$ on observe
\[
(S_t,A_t,R_{t+1},S_{t+1})=(s,a,r,s').
\]
Dans l'expression de $(T_*Q)(s,a)$, l'espérance
\[
\sum_{s',r} p(s',r\given s,a)\left[r+\gamma \max_{a'} Q(s',a')\right]
\]
est la moyenne de la variable aléatoire
\[
Y_t = R_{t+1}+\gamma \max_{a'} Q(S_{t+1},a')
\]
conditionnellement à $(S_t=s,A_t=a)$.

\begin{propriete}
Si les transitions sont générées par le MDP, alors
\[
\E[Y_t \given S_t=s,A_t=a] = (T_*Q)(s,a).
\]
\end{propriete}

\paragraph{Justification.}
C'est exactement la définition d'une espérance conditionnelle discrète :
\[
\E[Y_t \given s,a]
=
\sum_{s',r} p(s',r\given s,a)\left[r+\gamma\max_{a'}Q(s',a')\right].
\]

Le terme
\[
R_{t+1}+\gamma \max_{a'} Q(S_{t+1},a')
\]
est donc un \emph{estimateur non biaisé} de la cible de Bellman conditionnelle à l'état-action observé.

\section{Idée de correction incrémentale}

Si l'on connaissait $q_*$, on aurait idéalement
\[
Q(s,a)= (T_*Q)(s,a).
\]
Quand cette égalité n'est pas satisfaite, on mesure l'erreur locale par
\[
\delta_t = R_{t+1}+\gamma \max_{a'} Q(S_{t+1},a') - Q(S_t,A_t).
\]
Cette quantité est appelée \emph{erreur de Bellman échantillonnée} ou \emph{TD error}.

\begin{rappel}
Une correction incrémentale standard vers une cible $y_t$ s'écrit
\[
Q_{\text{new}} = Q_{\text{old}} + \alpha (y_t-Q_{\text{old}}),
\]
où $\alpha \in (0,1]$ est un pas d'apprentissage.
\end{rappel}

En prenant ici
\[
y_t = R_{t+1}+\gamma \max_{a'} Q(S_{t+1},a'),
\]
on obtient immédiatement la mise à jour de Q-Learning :
\[
Q(S_t,A_t)\leftarrow Q(S_t,A_t)+\alpha\left[R_{t+1}+\gamma \max_{a'}Q(S_{t+1},a')-Q(S_t,A_t)\right].
\]

\section{Pourquoi cette formule est naturelle}

\subsection{Réécriture convexe}

La mise à jour peut se réécrire
\[
Q_{\text{new}}(S_t,A_t)
=
(1-\alpha)\,Q_{\text{old}}(S_t,A_t)
+
\alpha \left[R_{t+1}+\gamma\max_{a'}Q_{\text{old}}(S_{t+1},a')\right].
\]

\paragraph{Interprétation.}
La nouvelle valeur est une moyenne pondérée entre :
\begin{itemize}[leftmargin=1.6em]
    \item l'ancienne estimation ;
    \item une cible issue de la transition observée.
\end{itemize}

\subsection{Pourquoi le maximum apparaît}

Dans l'équation de Bellman optimale, après être arrivé en $s'$, l'agent choisit l'action qui maximise la valeur future. C'est pour cela que la cible contient
\[
\max_{a'}Q(S_{t+1},a').
\]

\begin{remarque}
Cette maximisation définit une cible \emph{greedy}. Même si l'agent se comporte avec une politique exploratoire, la cible d'apprentissage correspond à la politique optimale recherchée.
\end{remarque}

\section{Caractère off-policy}

\begin{definition}[Off-policy]
Un algorithme est dit off-policy si la politique utilisée pour générer les données peut être différente de la politique implicitement évaluée ou améliorée.
\end{definition}

Q-Learning est off-policy car :
\begin{itemize}[leftmargin=1.6em]
    \item les données peuvent être générées par une politique $\epsilon$-greedy, ou toute autre politique suffisamment exploratoire ;
    \item la cible utilise malgré tout
    \[
    \max_{a'}Q(S_{t+1},a'),
    \]
    c'est-à-dire la politique greedy associée à $Q$.
\end{itemize}

\section{Convergence : idée générale}

Dans le cadre tabulaire fini, avec :
\begin{itemize}[leftmargin=1.6em]
    \item exploration suffisante ;
    \item pas d'apprentissage $(\alpha_t)$ vérifiant les conditions de Robbins-Monro
    \[
    \sum_t \alpha_t = \infty,
    \qquad
    \sum_t \alpha_t^2 < \infty ;
    \]
    \item récompenses bornées ;
\end{itemize}
Q-Learning converge vers $q_*$.

\paragraph{Idée mathématique.}
La mise à jour est une approximation stochastique du point fixe
\[
Q=T_*Q.
\]
Comme $T_*$ est une contraction pour la norme infinie dans le cadre actualisé $\gamma<1$, les corrections aléatoires convergent vers l'unique point fixe.

\section{Algorithme pas à pas}

\begin{enumerate}[leftmargin=1.8em]
    \item Initialiser $Q(s,a)$, par exemple à $0$.
    \item Choisir une politique de comportement, souvent $\epsilon$-greedy par rapport à $Q$.
    \item À chaque étape, observer $(S_t,A_t,R_{t+1},S_{t+1})$.
    \item Calculer la cible
    \[
    y_t = R_{t+1}+\gamma \max_{a'} Q(S_{t+1},a').
    \]
    \item Mettre à jour
    \[
    Q(S_t,A_t)\leftarrow Q(S_t,A_t)+\alpha(y_t-Q(S_t,A_t)).
    \]
    \item Répéter jusqu'à stabilisation de $Q$ ou jusqu'à un budget d'épisodes fixé.
\end{enumerate}

\section{Application à FrozenLake}

\subsection{Pourquoi FrozenLake convient}

FrozenLake possède :
\begin{itemize}[leftmargin=1.6em]
    \item un nombre fini d'états ;
    \item quatre actions discrètes ;
    \item des transitions stochastiques si \texttt{is\_slippery=True} ;
    \item une récompense rare : $1$ à l'arrivée, $0$ sinon.
\end{itemize}
C'est donc un exemple classique pour l'apprentissage tabulaire.

\subsection{Cas non glissant}

Quand \texttt{is\_slippery=False}, la transition est déterministe. La cible observée
\[
R_{t+1}+\gamma \max_{a'}Q(S_{t+1},a')
\]
est plus stable, donc l'apprentissage est plus rapide.

\subsection{Cas glissant}

Quand \texttt{is\_slippery=True}, la même action peut conduire à plusieurs états suivants. L'espérance exacte n'est pas calculée explicitement, mais les échantillons successifs finissent par l'approcher. C'est précisément le principe de Q-Learning.

\subsection{Lecture pédagogique}

Dans une interface FrozenLake, on peut :
\begin{itemize}[leftmargin=1.6em]
    \item afficher dans chaque case les quatre valeurs $Q(s,a)$ ;
    \item déduire la politique greedy
    \[
    \pi_{\text{greedy}}(s)=\argmax_a Q(s,a) ;
    \]
    \item comparer l'évolution de l'apprentissage entre la version glissante et la version non glissante.
\end{itemize}

\section{Résumé}

Q-Learning provient directement de l'équation d'optimalité de Bellman en remplaçant l'espérance par une transition échantillonnée et en conservant le maximum sur l'action future. Sa mise à jour
\[
Q(S_t,A_t)\leftarrow Q(S_t,A_t)+\alpha\left[R_{t+1}+\gamma \max_{a'}Q(S_{t+1},a')-Q(S_t,A_t)\right]
\]
est une correction incrémentale vers le point fixe optimal. Dans FrozenLake, il apprend progressivement quelles actions augmentent la probabilité d'atteindre la case but.

\end{document}
